%%-*-latex-*-

\documentclass[wide]{slides}

% Language
%
\usepackage{hyphenat}          % \hyp{} is a breakable dash
\usepackage{alltt}

\frenchspacing  % Follow French conventions after a period

% Graphics
%
\usepackage{graphicx}

% Maths & Logic
%
\usepackage{amsmath,amssymb,amsthm,stmaryrd}
\usepackage{mathpartir}

% ----------------------------------------------------------------
% Document
%

\maintitle{LIGO/Marigold: Smart Contracts and Layer-2}
\mainauthor{\textbf{Christian Rinderknecht}\\
  {\small\url{Christian.Rinderknecht@nomadic-labs.com}}\\
Nomadic Labs}
\confname{Tezos London Meetup}
\confshortname{Tezos}
\confdate{15 August 2019}

\begin{document}

\maketitle

\begin{slide}
  \title{A personal introduction}

  \begin{itemize}

    \item My alma mater is \emph{Universit\'e Pierre et Marie Curie}
      (UPMC, a.k.a. Paris~6).

    \item I did my doctoral studies at Inria, one of the most
      prestigious research institutes in informatics in France.

    \item I was a member of the team that developed the programming
      language OCaml.

    \item I went on to work as an engineer, a researcher and a
      professor for many years, across several countries (France,
      Korea, Hungary, Sweden), both in academia and private
      companies.

    \item In \oldstylenums{2018}, I joined Nomadic Labs, where most of
      the maintenance of the Tezos blockchain is done. My expertise is
      in compiler construction and functional programming. I have been
      working on a high-level language for writing smart contracts on
      Tezos.

  \end{itemize}

\end{slide}


\begin{slide}
  \title{A personal introduction}

  \begin{minipage}{0.5\linewidth}
    My book about functional programming is published in London!
  \end{minipage}%
  \begin{minipage}{0.5\linewidth}
    \includegraphics[scale=0.35]{my_book.eps}
  \end{minipage}%

\end{slide}

\begin{slide}
  \title{Nomadic Labs}

  \begin{itemize}

    \item Nomadic Labs is an engineering company based in Paris.

    \item It is funded by the Tezos Foundation, which tasked it with
      developing the Tezos ecosystem.

    \item It focuses on R\&D and is the main contributor to the Tezos
      blockchain itself.

    \item Most of the \oldstylenums{35}~engineers are doctors in
      informatics, specialised in programming language theory and
      practice, functional programming (OCaml), distributed systems,
      and formal verification.

    \item Nomadic Labs is also the home of Cortez, a mobile wallet for
      Tezos.

    \item It also helps with the training of Tezos users and other
      trainers.

  \end{itemize}

\end{slide}

\begin{slide}
  \title{Blockchains}

  \begin{itemize}

    \item The economic protocol of Tezos handles a ledger of
      transactions as a \textbf{blockchain}. In a nutshell: a
      replicated database with strong access control, immutable
      entries and resilience to malicious users.

    \item Those transactions consist in the exchange of assets
      (tokens) between peers of the network. Each peer has an address
      and an account (an amount of tokens), but \textbf{a physical
        person can be several peers.}

    \item \textbf{Public-key cryptography} is used to secure the
      identity of the senders of transactions and to ensure that past
      transactions are not tempered with. An address is the hash of
      the peer's public key.

    \item The business logic between the peers is \emph{not} recorded
      in the chain, that is, the knowledge of what they agree should
      trigger transfers of tokens.

    \item The economic protocol of Tezos offers the possibility to
      record in the chain that logic as a \textbf{smart contract}.

  \end{itemize}

\end{slide}

\begin{slide}
  \title{Smart Contracts on Tezos}

  \begin{itemize}

    \item A smart contract is a peer that also has a program
      associated to a private storage.

    \item That storage is writable only by the program, but is
      publicly readable (except by other contracts, to make them more
      deterministic by not depending on the current storage of another
      contract).

    \item The smart contract (by synecdoche) is executed each time a
      specific transaction is sent to its address.

    \item The transaction may include some parameters.

    \item Any call to a contract in a block is replicated by all the
      nodes in the chain, to check whether the block is valid before
      including it in their local view of the head of the chain.

    \item Once the block is validated, it is broadcasted by the
      underlying peer-to-peer, gossip network (as usual).

  \end{itemize}

\end{slide}

\begin{slide}
  \title{Smart Contracts on Tezos}

  \begin{itemize}

    \item The validation of smart contracts therefore adds a
      \textbf{delay} in the validation of a block: the execution time of
      all the smart contracts in it must fit the interblock time for
      the chain (about one minute on Tezos).

    \item Therefore, each smart contract is allowed a given and fixed
      quantity of computation, measured in \textbf{gas}.

    \item Each instruction has an associated cost expected to be
      proportional to the wall-clock time and an \textbf{ad hoc
        estimation} is given by the node's client.

    \item A \textbf{node} is a server process that is accessed by a
      command-line client or RPCs. A node comprises a view of the
      chain so far, the \textbf{context of the chain} (a map from
      addresses to token amounts, used to validate transactions), and
      the \textbf{mempool}.

    \item When the execution exceeds the alloted gas, it is stopped,
      its effect on the storage is rolled back, a failed transaction
      is included in the block, and thus fees are collected to
      \textbf{discourage spamming attacks}.

    \item The economic protocol sets limits on gas per block and per
      transaction.

  \end{itemize}

\end{slide}

\begin{slide}
  \title{Smart Contracts on Tezos}

  \begin{itemize}

    \item Each smart contract has a \textbf{code size} and a
      \textbf{storage size}, which are allocated on every node on the
      chain (in their context).

    \item To limit the memory needed to synchronise the chain and
      store it, a fee is set per byte and collected when a contract is
      deployed (\emph{originated}, in Tezos parlance).

    \item It is only at the normal termination of a contract that the
      effects of a contract becomes atomically visible to the network:
      the storage may appear modified and a list of operations may
      have been returned (transactions, contract creations and
      delegate settings) and validated.

    \item In particular, a smart contract can transfer tokens to other
      smart contracts, enabling the design of \textbf{complex
        distributed applications} in Tezos.

  \end{itemize}

\end{slide}

\begin{slide}
  \title{The design of Michelson}

  \begin{itemize}

    \item The language native to the Tezos blockchain for writing
      smart contracts is \textbf{Michelson}, a Domain-Specific
      Language (DSL) inspired by Lisp and Forth.

    \item This unusual lineage aims at satisfying unusual constraints,
      but entails some tensions in the design.

    \item First, to measure stepwise gas consumption,
      \textbf{Michelson is interpreted}.

    \item On the one hand, to assess gas usage per instruction,
      instructions should be simple, which points to \textbf{low-level
        features} (a RISC-like language).

    \item On the other hand, it was originally thought that users will
      want to write in Michelson instead of lowering a language to
      Michelson, because the gas cost would otherwise be harder to
      predict.

    \item This means that \textbf{high-level features} were deemed
      necessary (like a restricted variant of Lisp lambdas, a way to
      encode algebraic data types, as well as built-in sets, maps and
      lists).

  \end{itemize}

\end{slide}

\begin{slide}
  \title{The design of Michelson}

  \begin{itemize}

    \item To avoid ambiguous and otherwise misleading contracts, the
      layout of Michelson contracts has been constrained (e.g.,
      indentation, no UTF-8), and a \textbf{canonical form} was
      designed and enforced when storing contracts on the chain.

    \item To avoid the origination of contracts that would fail due to
      inconsistent assumptions on the storage and temporary values, a
      \textbf{strong static type system} was designed for Michelson.

    \item To reduce the size of the code, Michelson was designed as
      \textbf{a stack-based language}, whence the lineage from Forth
      and other concatenative languages like PostScript, Joy, Cat,
      Factor etc. (Java bytecode would count too.)

    \item Programs in those languages are \textbf{compact} because
      they assume an implicit stack in which some input values are
      popped, and output values are pushed, according to the current
      instruction being executed.

  \end{itemize}

\end{slide}

\begin{slide}
  \title{The semantics of Michelson}

  \begin{itemize}

    \item The semantics of each Michelson instruction is determined by
      the change of a prefix of the stack, that is, a segment starting
      at the top.

    \item An abstraction of that change can be captured by the
      \textbf{type of the instruction}.

    \item For example, the instruction \textsf{DUP} has type
      \(\alpha::A \rightarrow \alpha::\alpha::A\), meaning that if the
      stack before the instruction is executed has a topmost item
      (also called \emph{head}) of type~\(\alpha\), where the substack
      (also called \emph{tail}) has type~\(A\), then the instruction
      leaves a stack with two items of type~\(\alpha\), on top of a
      stack of type~\(A\). (We use the OCaml notation for consing,
      that is, pushing items.)

    \item This does not specify entirely the semantics, only the
      \textbf{static semantics}.

    \item We need a \textbf{dynamic semantics}, also called
      \emph{evaluation} or \emph{interpretation}, defined as a
      transition system, which states what happens to values in the
      stack. Here, the evaluation mimics exactly the typing judgement
      above:
\begin{mathpar}
\inferrule
    {\textsf{DUP} \; x::S \; \rightarrow \; x::x::S}
    {}
\end{mathpar}

  \end{itemize}

\end{slide}

\begin{slide}
  \title{The type system of Michelson}

  \begin{itemize}

    \item Whilst the types of instructions can be polymorphic, their
      instantiations must be monomorphic, hence instructions are not
      first-class values and cannot be partially interpreted.

    \item This enables a simple \textbf{static type checking}, as
      opposed to a complex type inference. It can be performed
      efficiently: \textbf{contract type checking consumes gas}.

    \item Basically, type checking aims at validating the composition
      of instructions, therefore is key to safely composing contracts
      (concatenation, activations).

    \item Once a contract passes type checking, it cannot fail due to
      inconsistent assumptions on the storage and other values (there
      are no null values, no casts), but it can still fail for other
      reasons: division by zero, token exhaustion, gas exhaustion, or
      an explicit \text{FAILWITH} instruction. This property is called
      \textbf{type safety}. Also, such a contract cannot remain stuck:
      this is the \textbf{progress} property.

    \item The existence of a formal type system for Michelson, of a
      formal specification of its dynamic semantics (evaluation), of a
      Michelson interpreter in Coq, of proofs in Coq of properties of
      some typical contracts, all those achievements are instances of
      \textbf{formal methods} in Tezos.

  \end{itemize}

\end{slide}

\begin{slide}
  \title{A voting contract in Michelson}

  \begin{itemize}

    \item The following example has been ripped off the paper
      \emph{Introduction to the Tezos Blockchain}, V. Allombert,
      M. Bourgoin, J. Tesson, HPCS, 2019, July 15–19, Dublin,
      Ireland.

    \item Disclaimer: They are my colleagues.

    \item We want to write a smart contract in Michelson that records
      votes of peers for their favourite candidate.

    \item For the sake of simplificity, a candidate is denoted by a
      string.

    \item The vote is open to all peers for a fee, and they have to
      pick one candidate in a fixed list.

    \item The state of the stack at a given point will be written in a
      comment starting with \texttt{\#}.

  \end{itemize}

\end{slide}

\begin{slide}
  \title{A voting contract in Michelson}

  \begin{itemize}

    \item We begin by defining the storage as a mapping from the
      candidates (\textsf{string}) to the number of corresponding
      votes (\textsf{int}):
      \begin{alltt}
        \textbf{storage} (map string int);
      \end{alltt}

    \item Then we specify the type of the parameter of the contract,
       that is, the vote:
      \begin{alltt}
        \textbf{parameter} string;
      \end{alltt}

    \item The execution starts with a stack containing a single pair
      made of the parameter and the storage:
      \begin{alltt}
        \textbf{code} \{
          # (vote, storage)
      \end{alltt}

  \end{itemize}

\end{slide}

\begin{slide}
  \title{A voting contract in Michelson}

  \begin{itemize}

    \item First, we check if the caller sent us enough tokens to
      vote, and, if not, we fail:
      \begin{alltt}
          AMOUNT;
          # amount::(vote, storage)
          PUSH mutez 5000;
          # 5000::amount::(vote, storage)
          IFCMPGT \{ PUSH string "Insufficient fee."; FAILWITH \} \{\};
          # (vote, storage)
      \end{alltt}

    \item The instruction \texttt{AMOUNT} pushes onto the stack the
      amount sent by the voter.

    \item The instruction \texttt{PUSH} pushes a typed constant.

    \item The instruction \texttt{IFCMPGT} compares the first two
      numbers on the stack, and, if the topmost is greater than the
      next, the first sequence (between braces) is executed, else the
      second one (which is empty here because of \texttt{FAILWITH}).

    \item The instruction \texttt{FAILWITH} reads the value on top of
      the stack, terminates the execution, and signals an error
      (denoted by the value) to the node and the client.

  \end{itemize}

\end{slide}

\begin{slide}
  \title{A voting contract in Michelson}

  \begin{itemize}

    \item In order to check the current number of votes for the same
      candidate, we start by duplicating the topmost item:
      \begin{alltt}
          DUP;
          # (vote, storage)::(vote, storage)
      \end{alltt}

    \item We then destructure the first pair:
      \begin{alltt}
          UNPAIR;
          # vote::storage::(vote, storage)
      \end{alltt}

    \item We can now consult the map in storage for the current count:
      \begin{alltt}
          GET;
          # (Some current | None)::(vote, storage)
      \end{alltt}
      The top of the stack is \texttt{None} if \texttt{vote} was not a
      valid choice, and \texttt{(Some current)} if there were
      \texttt{current} votes already.
  \end{itemize}

\end{slide}

\begin{slide}
  \title{A voting contract in Michelson}

  \begin{itemize}

    \item If \texttt{None}, we must fail; otherwise, we increment the
      number of votes:
      \begin{alltt}
          IF\_SOME \{ PUSH int 1; ADD; SOME \}
                   \{ PUSH string "Unknown candidate."; FAILWITH \}
          # Some (current+1) :: (vote, storage)
      \end{alltt}

    \item We need now to update the storage with the new vote. For
      that, we need to fish out the vote:
      \begin{alltt}
          DIP \{ UNPAIR \};
          # Some (current+1) :: vote :: storage
      \end{alltt}
      The instruction \texttt{DIP} applies its sequence of
      instructions (here \texttt{UNPAIR}) to the item just below the
      top.

    \item Then we need to reorder the two first items, like so:
      \begin{alltt}
          SWAP;
          # vote :: Some (current+1) :: storage
      \end{alltt}

  \end{itemize}

\end{slide}

\begin{slide}
  \title{A voting contract in Michelson}

  \begin{itemize}

    \item We can now update the map:
      \begin{alltt}
          UPDATE;
          # storage'
      \end{alltt}

    \item All contracts must normally end by leaving in the stack a
      pair made of a list of operations and the new storage.

    \item In this instance, the list is empty because all we do is
      update the internal storage of the contract. So we push an empty
      list of operations and pair it with the new storage:
      \begin{alltt}
          NIL operation;
          PAIR;
          # ([], storage')
        \}
      \end{alltt}

  \end{itemize}

\end{slide}

\begin{slide}
  \title{The full Michelson contract}
  \begin{alltt}
\textbf{storage} (map string int);
\textbf{parameter} string;
\textbf{code} \{
  AMOUNT;
  PUSH mutez 5000;
  IFCMPGT \{ PUSH string "Insufficient fee."; FAILWITH \} \{\};
  DUP;
  UNPAIR;
  GET;
  IF\_SOME \{ PUSH int 1; ADD; SOME \}
            \{ PUSH string "Unknown candidate."; FAILWITH \}
  DIP \{ UNPAIR \};
  SWAP;
  UPDATE;
  NIL operation;
  PAIR;
\}
  \end{alltt}

\end{slide}

\begin{slide}
  \title{Assessment of the contract}

  \begin{itemize}

    \item We could improve upon this design by enforcing a deadline on
      the ballot, or by granting the right to add new candidates in
      exchange of a fee.

    \item Anyhow, we had to modify, sometimes deeply, the stack in
      order to fetch and prepare the data for the instructions.

    \item Moreover, some instructions above were actually
      \textbf{macros} (\texttt{IF\_SOME}, \texttt{IFCMPGT}), which are
      actually expanded into a sequence of atomic instructions by the
      Michelson parser, so the canonical contract on the chain is
      longer.

    \item Programming in \textbf{Michelson is fun} and might even save
      the Caps Lock key from extinction.

    \item \emph{But does it scale up to larger and involved contracts
      with high stakes?}

    \item This is where \textbf{LIGO comes into play}.

  \end{itemize}

\end{slide}

\begin{slide}
  \title{LIGO}

  \begin{itemize}

    \item Like Michelson, \textbf{LIGO is a programming language for
      writing smart contracts on Tezos.}

    \item Unlike Michelson, LIGO is a language akin to what a
      mainstream programmer would expect.

    \item This means that LIGO features variables, expressions,
      function calls, data types, pattern matching etc.

    \item Nevertheless LIGO remains a DSL for the Tezos blockchain, so
      not all the usual bells and whistles of a mainstream language
      are available.

    \item It is hosted here: \url{https://ligolang.org/}

    \item It is an \textbf{open source, collective project} (under the
      MIT licence).

    \item Go and \textbf{try the tutorial}!

    \item (If you guessed why the name ``Michelson'', you will guess
      why ``LIGO''.)

  \end{itemize}

\end{slide}

\begin{slide}
  \title{LIGO}

  \begin{itemize}

    \item Perhaps the most striking feature of LIGO is that it comes
      in \textbf{different concrete syntaxes}, and even
      \textbf{different programming paradigms}.

    \item In other words, LIGO is not defined by one syntax and one
      paradigm, like imperative versus functional.

    \item There is \textbf{PascaLIGO}, which is inspired by Pascal,
      hence is an imperative language with lots of keywords, where
      values can be locally mutated after they have been annotated
      with their types (declaration).

    \item There is \textbf{CameLIGO}, which is inspired by the pure
      subset of OCaml, hence is a functional language with few
      keywords, where values cannot be mutated, but still require type
      annotations (unlike OCaml, whose compiler performs almost full
      type inference).

    \item \textbf{And more to come!} You can propose your own
      front-end for LIGO and we can help with its integration in the
      distribution.

  \end{itemize}

\end{slide}

\begin{slide}
  \title{A voting contract in PascaLIGO}

  \begin{alltt}
\textbf{type} store \textbf{is} map (string, nat)
\textbf{type} return  \textbf{is} list (operation) * store

\textbf{function} vote (\textbf{const} ballot : string; \textbf{var} store : store) : return \textbf{is}
  \textbf{begin}
    \textbf{if}   amount <= 5000mutez
    \textbf{then} failwith ("Insufficient fee.")
    \textbf{else case} store[ballot] \textbf{of}
                         None -> failwith ("Unknown candidate.")
          | Some (current) -> store[ballot] := current + 1n
          \textbf{end}
  \textbf{end with} ((\textbf{nil} : list (operation)), store)
  \end{alltt}

  \begin{itemize}

    \item \textbf{Easier to understand}, even though the generated
      Michelson is not so neat.

  \end{itemize}

\end{slide}

\begin{slide}
  \title{A voting contract in CameLIGO}

  \begin{alltt}
\textbf{type} store = (string, nat) map
\textbf{type} return = operation list * store

\textbf{let} vote (ballot, store : string * store) : return =
  \textbf{if}   amount <= 5000mutez
  \textbf{then} (failwith "Insufficient fee." : return)
  \textbf{else match} Map.find\_opt ballot store \textbf{with}
           None -> (failwith "Unknown candidate." : return)
        | Some current ->
             \textbf{let} store = Map.update ballot (Some (current + 1n)) store
             \textbf{in} (\textbf{nil} : operation list), store
  \end{alltt}

\end{slide}

\begin{slide}
  \title{A crowdfunding contract in PascaLIGO}

  \begin{itemize}

    \item A simple crowdfunding contract can be called in three
      manners:
      \begin{enumerate}

        \item a backer sends some funds before the deadline has
          passed, and only once;

        \item a backer claims their funds after the deadline has
          passed and the goal has not been reached;

        \item the owner withdraws the funds after the deadline has
          passed and the goal has been reached;

        \item all other cases are errors.

      \end{enumerate}

      \item First, we define the storage (or \emph{store}):
  \begin{alltt}
\textbf{type} store \textbf{is} \textbf{record}
    goal      : tez;
    deadline : timestamp;
    backers  : map (address, tez);
    funded   : bool;                  // Optional terminator semicolon
  \textbf{end}
\textbf{type} return \textbf{is} list (operation) * store
  \end{alltt}

  \end{itemize}

\end{slide}

\begin{slide}
  \title{A crowdfunding contract in PascaLIGO}

  \begin{itemize}

  \item The function to contribute to the crowdfunding:
    \smallskip
    \begin{alltt}
\textbf{function} back (\textbf{var} param : unit; \textbf{var} store : store) : return \textbf{is}
  \textbf{begin}
    \textbf{if}   now > store.deadline
    \textbf{then} failwith ("Past deadline.")
    \textbf{else} \textbf{case} store.backers[sender] \textbf{of}
             None -> store.backers[sender] := amount
          | Some (x) -> \textbf{skip}
          \textbf{end}
  \textbf{end with} ((\textbf{nil}: list (operation)), store)
    \end{alltt}

  \end{itemize}

\end{slide}

\begin{slide}
  \title{A crowdfunding contract in PascaLIGO}

  \begin{itemize}

    \item The function to claim funds is as follows:
    \smallskip
      \begin{alltt}
\textbf{function} claim (\textbf{var} param : unit; \textbf{var} store : store) : return \textbf{is}
  \textbf{begin}
    \textbf{var} op : list (operation) := \textbf{nil};
    \textbf{if} now <= store.deadline \textbf{then} failwith ("Too early.") \textbf{else}
    \textbf{case} store.backers[sender] \textbf{of}
        None -> failwith ("Not a backer.")
    | Some (asset) ->
          \textbf{if} balance >= store.goal \textbf{or} store.funded
          \textbf{then} failwith ("Goal reached: no refund.")
          \textbf{else} \textbf{begin}
                  \textbf{const} dest : contract (unit) = get_contract (sender);
                  op := \textbf{list} transaction (Unit, asset, dest) \textbf{end};
                  \textbf{remove} sender \textbf{from map} store.backers
                \textbf{end}
     \textbf{end}
  \textbf{end with} (op, store)
      \end{alltt}
  \end{itemize}

\end{slide}

\begin{slide}
  \title{A crowdfunding contract in PascaLIGO}

  \begin{itemize}

    \item The function to withdraw the funds is as follows:
      \smallskip
      \begin{alltt}
\textbf{function} withdraw (\textbf{var} param : unit; \textbf{var} store : store) : return \textbf{is}
  \textbf{begin}
    \textbf{var} op : list (operation) := \textbf{nil};
    \textbf{if} sender = owner \textbf{then}
      \textbf{if} now >= store.deadline \textbf{then}
        \textbf{if} balance >= store.goal \textbf{then}
           \textbf{begin}
             store.funded := True;
             op := \textbf{list} transaction (unit, owner, balance)] \textbf{end};
           \textbf{end}
        \textbf{else} failwith ("Below target.")
      \textbf{else} failwith ("Too early.")
    \textbf{else skip}
  \textbf{end with} (op, store)
\end{alltt}

  \end{itemize}

\end{slide}

\begin{slide}
  \title{Some peculiarities of LIGO}

  \begin{itemize}

    \item LIGO features multiple syntaxes and paradigms.

    \item Even within PascaLIGO, two styles are possible: terse or
      verbose. We illustrated the verbose style here. We plan to offer
      automatic style checking and two-way conversion by
      pretty\hyp{}printing.

    \item LIGO features \textbf{high-order quotations}, but not
      full-fledged closures.

    \item This means that, at the call site, \textbf{the arguments
      and the environment are always copied}, therefore any mutation
      (in PascaLIGO) will have no effect on the caller's arguments or
      the environment, except with built-in iterators.

    \item Quotations can be passed as arguments to others.

  \end{itemize}

\end{slide}

\begin{slide}
  \title{Quotations in LIGO}

\begin{alltt}
\textbf{function} iter (\textbf{const} delta : int; \textbf{const} l : list (int)) : int \textbf{is}
  \textbf{var} acc : int := 0;
  \textbf{procedure} aggregate (\textbf{const} i : int) \textbf{is}
    \textbf{begin}
      acc := acc + i
    \textbf{end}
  \textbf{begin}
    aggregate (delta);         // Has no effect on acc
    list\_iter (l, aggregate); // Has an effect on acc
  \textbf{end with} acc
\end{alltt}

\end{slide}

\begin{slide}
  \title{Tooling for LIGO}

  \begin{itemize}

    \item Several tools are currently being developed, aiming at
      facilitating the adoption of LIGO.

    \item \textbf{An IDE based on VSCode} has been made available,
      featuring
      \begin{itemize}

        \item syntax highlighting,

        \item one-click compilation to Michelson,

        \item \emph{dry runs} to locally execute contracts on a
          sandbox,

        \item a reactive counter estimating the gas consumption.

      \end{itemize}

    \item The architecture of VSCode, with its Language Server
      Protocol, hopefully opens the door to plugins written in any
      programming language, e.g., static analysis in OCaml.

    \item \textbf{A web-based IDE} with the same set of features.

    \item \textbf{An SDK} to make it simple for developers to make
      applications involving LIGO contracts, e.g., mobile smart
      wallets or web applications.

  \end{itemize}

\end{slide}

\begin{slide}
  \title{Research \& Development on the LIGO compiler}

  \begin{itemize}

    \item We are working on a more \textbf{powerful type system} which
      will enable the writing of more expressive contracts, featuring
      more type inference (less annotations) and enabling a greater
      variety of programming paradigms (e.g., object-oriented).

    %% \item We are working on a \textbf{generic front-end} aimed at
    %%   replacing the different front-ends (lexers and parsers) we
    %%   currently have. We also experiment with in-contract syntax
    %%   extensions.

    \item We are writing a \textbf{certified backend in Coq}, that is,
      a Michelson code generator proven correct and extracted to OCaml
      from its specification.

    \item Those endeavours are not just engineering, they are
      instances of \textbf{applied research} and require a strong
      background on programming language theory.

    \item But... you are welcome to join our effort, research or
      development, even if you do not have a PhD, of course.

  \end{itemize}

\end{slide}

\begin{slide}
  \title{The team behind LIGO}
  \begin{itemize}

    \item I wrote ``we'' and I introduced myself as working for
      Nomadic Labs.

    \item Nomadic Labs has been helping with spinning off LIGO into a
      company, which I will formally join next month. Both companies
      are funded by the Tezos Foundation.

    \item \textbf{We are an international team.}

    \item Five of us are based in Paris, and eight are contractors
      spread around the globe, all working as one team (it is not just
      about deliverables).

    \item I became the COO, but I also coach new hires and work on an
      experimental parser.

    \item If you are passionate about compilers, distributed
      applications on the blockchain, the web, please get in touch
      with us.

  \end{itemize}

\end{slide}

\begin{slide}
  \title{Marigold}

  \begin{itemize}

    \item The LIGO project is the first stage of a larger project
      called \textbf{Marigold}.

    \item Marigold aims at \textbf{accelerating transactions},
      including smart contracts, with \textbf{layer-2 innovations}
      (sometimes not adequately called \emph{off-chain} or
      \emph{side-chain} solutions).

    \item We plan to amend the Tezos protocol itself to create
      \textbf{tokens (Non Fungible Tokens) with Plasma}, which could
      be done with smart contracts, but much slower.

    \item This is possible because Tezos features a
      \textbf{metaconsensus protocol}, which empowers peers to propose
      amendments (i.e., source code) and submit their patches to a
      vote.

    \item For LIGO, we already proposed some Michelson instructions to
      make the code generated by the LIGO compiler faster (fetching
      and committing data deep within the stack).

    \item Cryptoeconomics Labs (Japan)
      (\url{https://www.cryptoeconomicslab.com}) work on a
      \textbf{general Plasma}, that is, a Plasma that is extended to
      encompass a larger class of contracts than just tokens, and of
      interest to Marigold.

  \end{itemize}

\end{slide}

\end{document}
